%% kapa: exam=1 round=28 year=2017 period=1 form=A subject=real_estate subject_ko=부동산학원론 seq=24 number=104
%% source: Q-Net 감정평가사 기출문제 (cst00309) artlSeq=5199603 file=extracted/1008/5199603/2017년 제28회 감정평가사 1차시험 문제지/1교시 A형.pdf
%% conversion: manual transcription from page image; review=proofread-gpt-terra; flags=none
\begin{question}{104}
환원이율에 관한 설명으로 옳지 않은 것은?
\begin{choices}
  \item 환원이율은 투하자본에 대한 수익비율로써 상각 후ㆍ세공제전의 이율을 말한다.
  \item 개별환원이율이란 토지와 건물 각각의 환원이율을 말한다.
  \item 환원이율이란 대상부동산이 장래 산출할 것으로 기대되는 표준적인 순수익과 부동산가격의 비율이다.
  \item 환원이율은 순수익을 자본환원해서 수익가격을 구하는 경우에 적용되며, 이는 결국 부동산의 수익성을 나타낸다.
  \item 세공제전 환원이율이란 세금으로 인한 수익의 변동을 환원이율에 반영하여 조정(배제)하지 않은 환원이율을 말한다.
\end{choices}
\end{question}
\answer{1}
